\par\bigskip
\SourceLeaf{421}
\begin{center}
\pdfbookmark[2]{Gauss to Olbers, 26 November 1810}{pallas-brief-18101126}
\textsc{Gauss} to \textsc{Olbers}, 26 November 1810.
\end{center}
\noindent\hspace*{\parindent}I believe I already wrote to you in my last letter that I intended to undertake a study\HistoricalNoteEight{See p. 473 of this volume.]} of the perturbations of Pallas by Jupiter. I meant to work out their magnitude for the years from 1803 to 1811 alone, leaving aside their general theory. The calculation of three sets of elliptic elements from the oppositions of 1803, 1804, 1805, 1807 — 1804, 1805, 1807, 1808 — 1805, 1807, 1808, 1809 showed how strong the effect of the perturbations was. Judge for yourself from the following samples:
\begin{center}
\begingroup\footnotesize\setlength{\tabcolsep}{3pt}\renewcommand{\arraystretch}{1.2}
\begin{tabular}{lccc}
Perihelion 1803\ .\,.\,.\,.\,.\,.\,.\,. & \(121^{0}\,3'11''\) & \(121^{0}\,5'22''\) & \(120^{0}58'\,5''\)\\
\(\varphi\)\ .\,.\,.\,.\,.\,.\,.\,.\,.\,.\,.\,.\,.\,.\,.\,.\,.\,.\,.\,.\,. & \(14^{0}10'59''\) & \(14^{0}10'\,4''\) & \(14^{0}\,9'37''\)\\
Mean daily motion.\ .\,.\,. & \(770''{,}214\) & \(770''{,}447\) & \(770''{,}926.\)\\
\end{tabular}\endgroup
\end{center}
\par
\noindent\hspace*{\parindent}I became curious whether these six oppositions could be reconciled much better if Jupiter's perturbations were taken into account. I confess that work of this kind, unpleasant in itself, had held too little interest for me as long as a single ellipse represented all the observations, and even the prize of 6000 francs could hardly have tempted me to it. I have now taken heart and set about the work. After devoting most of my available time to it for 4 weeks, I have finished the main part, so that, unless my calculations contain major errors, I may hope soon to reap the fruits. I thought it most convenient to treat the elements as variable and to compute their daily changes due to Jupiter at intervals of 50 days for a period of more than 8 years. I began by developing formulas of my own, or rather by putting those I had developed several years ago into the most manageable form I could. I found \textsc{Laplace}'s far too cumbersome and would certainly have needed two or three times as long with his as with mine. I calculated Jupiter's positions from \textsc{Bouvard}'s tables (using a particular device\HistoricalNoteDouble{See p. 475 of this volume.]}) \SourceLeaf{422}and the positions of Pallas from a 4th set of elements fitted as closely as possible to all the oppositions. Today I have at last finished those 61 daily changes for each of the six elements. There remain the mechanical integration (little more than addition) and the fresh calculation of the fundamental elements which, combined with the resulting perturbations, best represent the oppositions. These perturbations are much more considerable than I thought, and particularly strong in 1809, when Jupiter was in heliocentric conjunction with Pallas. I therefore expect my ephemeris to be considerably off at the next opposition. \textsc{Harding} recently looked for Pallas but is still unsure whether he found it. If you observe it, I would be much obliged if you would let me know promptly. The daily motion will now be more than \(2''\) greater than in 1803. The angle \(\varphi\), which has steadily decreased, will probably be \(13'\) smaller, and the perihelion has probably also receded by \(20'\) or \(30'\), etc. If the agreement with the six previous oppositions and the 7th next year proves satisfactory to me, I may also undertake the general theory, although the immense amount of mechanical computation it entails is very discouraging. Jupiter's perturbations are so predominant that those of the other planets may be left aside for the time being.\quad.\quad.\quad.\quad.\quad.
\par
\noindent\hspace*{\parindent}If, by putting in a word in Paris, you should bring about a further extension of the deadline for the prize problem, that would give me one more reason to undertake the work just mentioned. The success of my present work must, however, be the first condition. In a fortnight I shall know where I stand, at least with respect to the observations already made. I shall then be doubly curious about the forthcoming ones, which will first show the full strength of the perturbations suffered at the conjunction.\quad.\quad.\quad.\quad.\quad.
\par
\noindent\hspace*{\parindent}P. S. Before closing this letter I have also completed the integration of the perturbation of the inclination, which is the least considerable. If I set the inclination for the 1st opposition
\begin{center}in 1803\quad at\quad \(34^{0}37'47'',\)\end{center}
then it is
\SourceLeaf{423}
\begin{center}
\begin{tabular}{lc}
1804 & \(34^{0}37'36''\)\\
1805 & \(34^{0}37'12''\)\\
1807 & \(34^{0}37'25''\)\\
1808 & \(34^{0}37'31''\)\\
1809 & \(34^{0}36'40''\)\\
1811 & \(34^{0}35'\,5''.\)\\
\end{tabular}
\end{center}
\par
